\chapter*{Abstract}


North S\'{a}mi is an indigenous Uralic low-resource language, classified by UNESCO as \textit{definitely endangered} and faces many challenges in digitalization. The most accurate publicly available Optical Character Recognition (OCR) models for the language are large transformer based variants released by the National Library of Norway (NLN), their \trocrSize{}~MB size and over \trocrInferenceIntroSec{} seconds per image inference time on CPU make them unsuitable for the community driven heritage projects, mobile field documentation and unbudgeted institutions that would benefit most from them.

This bachelor thesis explores whether a much smaller architecture can approach the accuracy of the state of the art (SOTA) model closely enough to be useful. A modular OCR pipeline (5 backbones $\times$ 3 sequence encoders, with a shared CTC head) was built. Only four variants were benchmarked (due to resource constraints), alongside seven public Transformer OCR (TrOCR) variants, on a \testSamples{} sample test set covering both synthetic in domain text and historical out of domain content from Turi's \turiYear{} \textit{Muitalus s\'{a}miid birra}.

The headline contribution, \texttt{ctc\_simple}, is a 7 layer Convolutional Neural Network---Connectionist Temporal Classification (CNN-CTC) model trained from scratch. At \ctcSize{}~MB and \ctcTime{}~s per image on CPU, it reaches \ctcCER{}\% CER (Character Error Rate) (\ctcCharAcc{}\% character accuracy), within \cerGap{}\,pp in CER of the best transformer baseline \trocrBestCER{}\% CER (\trocrBestCharAcc{}\% character accuracy) while running roughly \speedupFactor{}$\times$ faster and being \sizeRatio{}$\times$ smaller. A linguistically motivated conjunction based sentence splitting step further reduces long line CER from \splitLongOriginalCER{}\% to \splitLongSplitCER{}\%, even if only as a proof of concept, and an end to end OCR$\rightarrow$Machine Translation (MT) pipeline against TartuNLP's North~S\'{a}mi English API demonstrates the model's usability for downstream tasks.

This summary report focuses on the project's context, goals, process and process reflections, providing more detailed background and a personal perspective on the work. The technical details, experiments and scientific claims are covered in the full conference paper accepted to IISA 2026, titled \textbf{Lightweight OCR in North S\'{a}mi Heritage Digitization for Smart Learning}, and in the extended abstract accepted to SCAI 2026, titled \textbf{OCR and Machine Translation Pipeline for Low Resource S\'{a}mi Languages}.